\chapter{2004年全国II卷（理）}
\section{单选题}
\begin{problem}
已知集合 \(M = \{x \mid x^2 < 4\}\)，\(N = \{x \mid x^2 - 2x - 3 < 0\}\)，则集合 \(M \cap N =\)（\(\quad\)）

\choices
    {\( \{x \mid x < -2\} \)}
    {\( \{x \mid x > 3\} \)}
    {\( \{x \mid -1 < x < 2\} \)}
    {\( \{x \mid 2 < x < 3\} \)}
\end{problem}

\begin{answer}
C
\end{answer}

\begin{solution}
由 \(x^2<4\) 得 \(-2<x<2\)，所以 \(M=(-2,2)\). 又因为
\(x^2-2x-3=(x+1)(x-3)\)，且二次函数开口向上，所以 \(N=(-1,3)\). 因此
\(M\cap N=(-1,2)=\{x\mid -1<x<2\}\)，故选 C.
\end{solution}

\begin{problem}
\(\displaystyle\lim_{x\to 1}\frac{x^2 + x - 2}{x^2 + 4x - 5} =\)（\(\quad\)）

\choices
    {\( \frac{1}{2} \)}
    {\(1\)}
    {\( \frac{2}{5} \)}
    {\( \frac{1}{4} \)}
\end{problem}

\begin{answer}
A
\end{answer}

\begin{solution}
当 \(x\ne1\) 时，\(x^2+x-2=(x-1)(x+2)\)，\(x^2+4x-5=(x-1)(x+5)\)，因而
\[
\displaystyle\lim_{x\to1}\frac{x^2+x-2}{x^2+4x-5}=\lim_{x\to1}\frac{x+2}{x+5}=\frac36=\frac12
\]
故选 A.
\end{solution}

\begin{problem}
设复数 \(\omega = -\frac{1}{2} + \frac{\sqrt{3}}{2}\i\)，则 \(1 + \omega =\)（\(\quad\)）

\choices
    {\(-\omega\)}
    {\(\omega^2\)}
    {\(-\frac{1}{\omega}\)}
    {\(\frac{1}{\omega^2}\)}
\end{problem}

\begin{answer}
C
\end{answer}

\begin{solution}
由 \(\omega=-\frac12+\frac{\sqrt3}{2}\i\)，有
\[
\omega^2=\left(-\frac12+\frac{\sqrt3}{2}\i\right)^2=-\frac12-\frac{\sqrt3}{2}\i
\]
且 \(\omega\omega^2=1\)，于是
\[
1+\omega=\frac12+\frac{\sqrt3}{2}\i=-\omega^2=-\frac1\omega
\]
故选 C.
\end{solution}

\begin{problem}
已知圆 \(C\) 与圆 \((x - 1)^2 + y^2 = 1\) 关于直线 \(y = -x\) 对称，则圆 \(C\) 的方程为（\(\quad\)）

\choices
    {\((x + 1)^2 + y^2 = 1\)}
    {\(x^{2} + y^{2} = 1\)}
    {\(x^{2} + (y + 1)^{2} = 1\)}
    {\(x^{2} + (y - 1)^{2} = 1\)}
\end{problem}

\begin{answer}
C
\end{answer}

\begin{solution}
圆 \((x-1)^2+y^2=1\) 的圆心为 \((1,0)\)，半径为 \(1\). 关于直线 \(y=-x\) 的对称变换为 \((x,y)\mapsto(-y,-x)\)，所以圆心变为 \((0,-1)\)，半径不变. 故圆 \(C\) 的方程为 \(x^2+(y+1)^2=1\)，选 C.
\end{solution}

\begin{problem}
已知函数 \( y = \tan (2x + \varphi) \) 的图象过点 \( \left(\frac{\pi}{12}, 0\right) \)，则 \( \varphi \) 可以是（\(\quad\)）

\choices
    {\(-\frac{\pi}{6}\)}
    {\(\frac{\pi}{6}\)}
    {\(-\frac{\pi}{12}\)}
    {\(\frac{\pi}{12}\)}
\end{problem}

\begin{answer}
A
\end{answer}

\begin{solution}
将点 \(\left(\frac\pi{12},0\right)\) 代入函数，得
\[
\tan\left(\frac\pi6+\varphi\right)=0
\]
因此 \(\frac\pi6+\varphi=k\pi\)，其中 \(k\in\Z\). 取 \(k=0\)，得 \(\varphi=-\frac\pi6\)，故选 A.
\end{solution}

\begin{problem}
函数 \( y = -\e^{x} \) 的图象（\(\quad\)）

\choices
    {与 \(y = \e^{x}\) 的图象关于 \(y\) 轴对称}
    {与 \( y = \e^{x} \) 的图象关于坐标原点对称}
    {与 \( y = \e^{-x} \) 的图象关于 \( y \) 轴对称}
    {与 \( y = \e^{-x} \) 的图象关于坐标原点对称}
\end{problem}

\begin{answer}
D
\end{answer}

\begin{solution}
函数 \(y=\e^{-x}\) 的图象上任一点为 \((t,\e^{-t})\). 关于坐标原点对称后，该点变为 \((-t,-\e^{-t})\). 令 \(x=-t\)，则纵坐标为 \(-\e^x\)，所以所得图象为 \(y=-\e^x\). 故 \(y=-\e^x\) 与 \(y=\e^{-x}\) 的图象关于坐标原点对称，选 D.
\end{solution}

\begin{problem}
已知球 \(O\) 的半径为 \(1\)，\(A\)，\(B\)，\(C\) 三点都在球面上，且每两点间的球面距离为 \(\frac{\pi}{2}\)，则球心 \(O\) 到平面 \(ABC\) 的距离为（\(\quad\)）

\choices
    {\(\frac{1}{3}\)}
    {\(\frac{\sqrt{3}}{3}\)}
    {\(\frac{2}{3}\)}
    {\(\frac{\sqrt{6}}{3}\)}
\end{problem}

\begin{answer}
B
\end{answer}

\begin{solution}
球面距离为 \(\frac\pi2\)，半径为 \(1\)，所以相应的圆心角为 \(\frac\pi2\)，从而
\(\overrightarrow{OA}\perp\overrightarrow{OB}\)、\(\overrightarrow{OB}\perp\overrightarrow{OC}\)、\(\overrightarrow{OC}\perp\overrightarrow{OA}\). 因此
\(AB^2=BC^2=CA^2=2\)，三角形 \(ABC\) 是边长为 \(\sqrt2\) 的正三角形.

设正三角形 \(ABC\) 的中心为 \(H\). 由于
\(\overrightarrow{OH}=\frac13(\overrightarrow{OA}+\overrightarrow{OB}+\overrightarrow{OC})\)，且
\[
\overrightarrow{OH}\cdot\overrightarrow{AB}
=\frac13(\overrightarrow{OA}+\overrightarrow{OB}+\overrightarrow{OC})\cdot(\overrightarrow{OB}-\overrightarrow{OA})=0
\]
又有
\(\overrightarrow{OH}\cdot\overrightarrow{AC}=\frac13(\overrightarrow{OA}+\overrightarrow{OB}+\overrightarrow{OC})\cdot(\overrightarrow{OC}-\overrightarrow{OA})=0\)，所以 \(OH\perp\) 平面 \(ABC\). 正三角形的外接圆半径为
\[
AH=\frac{\sqrt3}{3}\cdot\sqrt2=\frac{\sqrt6}{3}
\]
在直角三角形 \(OHA\) 中，
\[
OH=\sqrt{OA^2-AH^2}=\sqrt{1-\frac23}=\frac{\sqrt3}{3}
\]
故选 B.
\end{solution}

\begin{problem}
在坐标平面内，与点 \(A(1,2)\) 距离为 \(1\)，且与点 \(B(3,1)\) 距离为 \(2\) 的直线共有（\(\quad\)）

\choices
    {\(1\text{ 条}\)}
    {\(2\text{ 条}\)}
    {\(3\text{ 条}\)}
    {\(4\text{ 条}\)}
\end{problem}

\begin{answer}
B
\end{answer}

\begin{solution}
设所求直线为 \(l\)，取单位法向量 \(\bs n\)，并写成 \(\bs n\cdot X=c\). 由点到直线的距离条件，存在
\(\varepsilon_A\)，\(\varepsilon_B\in\{1,-1\}\)，使得
\(\bs n\cdot A-c=\varepsilon_A\)，\(\bs n\cdot B-c=2\varepsilon_B\)
相减得
\[
\bs n\cdot\overrightarrow{AB}=2\varepsilon_B-\varepsilon_A
\]
其中 \(\overrightarrow{AB}=(2,-1)\)，\(\lvert AB\rvert=\sqrt5\). 右端为 \(\pm3\) 时不可能，因为
\(\lvert\bs n\cdot\overrightarrow{AB}\rvert\leqslant\lvert AB\rvert=\sqrt5<3\).

当右端为 \(1\) 时，设 \(\bs n=(u,v)\)，则
\(2u-v=1\)，且 \(u^2+v^2=1\). 由 \(v=2u-1\) 得
\(u^2+(2u-1)^2=1\)，即 \(u(5u-4)=0\)，所以
\[
(u,v)=(0,-1)\quad\text{或}\quad (u,v)=\left(\frac45,\frac35\right)
\]
这两组单位法向量分别给出直线 \(y=3\) 和 \(4x+3y=5\). 右端为 \(-1\) 时得到的是这两条直线的相反法向量表示，不产生新直线. 因此共有两条直线，选 B.
\end{solution}

\begin{problem}
已知平面上直线 \(L\) 的方向向量 \(\bs{e} = \left(-\frac{4}{5}, \frac{3}{5}\right)\)，点 \(O(0,0)\) 和 \(A(1, -2)\) 在 \(L\) 上的射影分别是 \(O_{1}\) 和 \(A_{1}\)，则 \(\overrightarrow{O_{1}A_{1}} = \lambda \bs{e}\)，其中 \(\lambda =\)（\(\quad\)）

\choices
    {\(\frac{11}{5}\)}
    {\(-\frac{11}{5}\)}
    {\(2\)}
    {\(-2\)}
\end{problem}

\begin{answer}
D
\end{answer}

\begin{solution}
因为 \(\bs e\) 是单位向量，正交投影后两点的连线向量等于
\(\overrightarrow{OA}\) 在 \(\bs e\) 方向上的投影，因此
\[
\lambda=\overrightarrow{OA}\cdot\bs e=(1,-2)\cdot\left(-\frac45,\frac35\right)=-\frac45-\frac65=-2
\]
故选 D.
\end{solution}

\begin{problem}
函数 \( y = x\cos x - \sin x \) 在下面哪个区间内是增函数（\(\quad\)）

\choices
    {\(\left(\frac{\pi}{2}, \frac{3\pi}{2}\right)\)}
    {\((\pi, 2\pi)\)}
    {\(\left(\frac{3\pi}{2}, \frac{5\pi}{2}\right)\)}
    {\((2\pi, 3\pi)\)}
\end{problem}

\begin{answer}
B
\end{answer}

\begin{solution}
设 \(f(x)=x\cos x-\sin x\)，则 \(f'(x)=-x\sin x\). 四个选项中的区间都在正半轴上，所以在这些区间内，\(f'(x)>0\) 等价于 \(\sin x<0\).

在 \((\pi,2\pi)\) 内始终有 \(\sin x<0\)，所以 \(f\) 在该区间内递增. 其余三个区间都包含 \(\sin x>0\) 的子区间，函数在这些子区间上递减，不能在整个区间内递增. 因此选 B.
\end{solution}

\begin{problem}
函数 \( y = \sin^4 x + \cos^2 x \) 的最小正周期为（\(\quad\)）

\choices
    {\( \frac{\pi}{4} \)}
    {\( \frac{\pi}{2} \)}
    {\(\pi\)}
    {\(2\pi\)}
\end{problem}

\begin{answer}
B
\end{answer}

\begin{solution}
利用恒等式 \(\sin^2x+\cos^2x=1\)，可得
\[
f(x)=\sin^4x+\cos^2x=\frac78+\frac18\cos4x
\]
所以 \(\frac\pi2\) 是函数的一个正周期. 若正数 \(T\) 是函数的周期，则由 \(f(0)=f(T)\) 得
\[
\frac78+\frac18\cos4T=1
\]
从而 \(\cos4T=1\)，即 \(T=k\frac\pi2\)，其中 \(k\) 为正整数. 故最小正周期为 \(\frac\pi2\)，选 B.
\end{solution}

\begin{problem}
在由数字 \(1\)，\(2\)，\(3\)，\(4\)，\(5\) 组成的所有没有重复数字的 \(5\) 位数中，大于 \(23145\) 且小于 \(43521\) 的数共有（\(\quad\)）

\choices
    {\(56\text{ 个}\)}
    {\(57\text{ 个}\)}
    {\(58\text{ 个}\)}
    {\(60\text{ 个}\)}
\end{problem}

\begin{answer}
C
\end{answer}

\begin{solution}
按首位分类. 首位为 \(3\) 时，后四位任意排列，有 \(4!=24\) 个.

首位为 \(2\) 时，只有次位为 \(3\)、\(4\) 或 \(5\) 才可能大于 \(23145\). 次位为 \(3\) 时，后三位由 \(1\)，\(4\)，\(5\) 排列，除去等于 \(145\) 的一种，有 \(3!-1=5\) 个；次位为 \(4\) 或 \(5\) 时均自动满足，各有 \(3!=6\) 个，所以共有 \(5+6+6=17\) 个.

首位为 \(4\) 时，只有次位为 \(1\)、\(2\) 或 \(3\) 才可能小于 \(43521\). 次位为 \(1\) 或 \(2\) 时各有 \(3!=6\) 个；次位为 \(3\) 时，后三位由 \(1\)，\(2\)，\(5\) 排列，除去等于 \(521\) 的一种，有 \(3!-1=5\) 个，所以共有 \(17\) 个. 首位为 \(1\) 的数都小于 \(23145\)，首位为 \(5\) 的数都大于 \(43521\).

总数为 \(24+17+17=58\)，故选 C.
\end{solution}

\section{填空题}
\begin{problem}
从装有 \(3\text{ 个}\) 红球，\(2\text{ 个}\) 白球的袋中随机取出 \(2\text{ 个}\) 球，设其中有 \(\xi\text{ 个}\) 红球，则随机变量 \(\xi\) 的概率分布为：
\begin{center}
\begin{tabular}{|c|c|c|c|}
\hline
\(\xi\) & \(0\) & \(1\) & \(2\) \\
\hline
\(P\) &\fillinblank{} &\fillinblank{} &\fillinblank{} \\
\hline
\end{tabular}
\end{center}
\end{problem}

\begin{answer}
\(0.1\)，\(0.6\)，\(0.3\)
\end{answer}

\begin{solution}
从 \(5\) 个球中取出 \(2\) 个球的总数为 \(\mathrm C_5^2=10\). 没有红球、恰有一个红球、恰有两个红球的取法数分别为
\(\mathrm C_2^2=1\)、\(\mathrm C_3^1\mathrm C_2^1=6\)、\(\mathrm C_3^2=3\). 所以
\[
P(\xi=0)=\frac1{10}=0.1
\]
\(P(\xi=1)=\frac6{10}=0.6\)，\(P(\xi=2)=\frac3{10}=0.3\)，填空依次为 \(0.1\)、\(0.6\)、\(0.3\).
\end{solution}

\begin{problem}
设 \(x\)，\(y\) 满足约束条件 \(\begin{cases}
x \geqslant 0,\\
x \geqslant y,\\
2x - y \leqslant 1,
\end{cases}\) 则 \(z = 3x + 2y\) 的最大值是\fillinblank{}.
\end{problem}

\begin{answer}
5
\end{answer}

\begin{solution}
由 \(y\leqslant x\) 和 \(2x-y\leqslant1\) 得
\(x\leqslant2x-y\leqslant1\). 再结合 \(x\geqslant0\)，可知 \(0\leqslant x\leqslant1\). 又因为 \(y\leqslant x\)，所以
\[
z=3x+2y\leqslant5x\leqslant5
\]
当 \(x=y=1\) 时，三个约束条件均满足，且 \(z=5\)，故最大值为 \(5\).
\end{solution}

\begin{problem}
设中心在原点的椭圆与双曲线 \(2x^{2} - 2y^{2} = 1\) 有公共的焦点，且它们的离心率互为倒数，则该椭圆的方程是\fillinblank{}.
\end{problem}

\begin{answer}
\(\frac{x^{2}}{2} + y^{2} = 1\)
\end{answer}

\begin{solution}
双曲线可写为
\(\frac{x^2}{1/2}-\frac{y^2}{1/2}=1\)，所以其半实轴平方和半虚轴平方分别为 \(\frac12\) 和 \(\frac12\). 由双曲线焦半距公式，
\[
c^2=\frac12+\frac12=1
\]
故 \(c=1\)，离心率为 \(e_h=\frac{c}{1/\sqrt2}=\sqrt2\).

椭圆与双曲线有公共焦点，所以椭圆的焦半距也是 \(c=1\). 由离心率互为倒数，椭圆离心率为 \(e_e=1/\sqrt2\). 设椭圆长半轴为 \(a\)，则
\[
\frac ca=\frac1{\sqrt2}
\]
从而 \(a^2=2\). 再由 \(b^2=a^2-c^2\) 得 \(b^2=1\)，所以椭圆方程为 \(\frac{x^2}{2}+y^2=1\).
\end{solution}

\begin{problem}
下面是关于四棱柱的四个命题：

\begin{circlelist}
\item 若有两个侧面垂直于底面，则该四棱柱为直四棱柱；
\item 若两个过相对侧棱的截面都垂直于底面，则该四棱柱为直四棱柱；
\item 若四个侧面两两全等，则该四棱柱为直四棱柱；
\item 若四棱柱的四条对角线两两相等，则该四棱柱为直四棱柱.
\end{circlelist}

其中，真命题的编号是\fillinblank{}.（写出所有真命题的编号）
\end{problem}

\begin{answer}
\(\circled{2}\circled{4}\)
\end{answer}

\begin{solution}
对命题 \(\circled{1}\)，取底面为水平正方形，取侧棱方向向量为 \(\bs v=(1,0,1)\)，并取沿 \(x\) 轴方向的两条相对底边. 这两个侧面分别由 \((1,0,0)\) 与 \(\bs v\) 张成，均含有竖直方向 \((0,0,1)=\bs v-(1,0,0)\)，因而都是竖直平面，均垂直于底面；但 \(\bs v\) 在底面内有非零投影，侧棱不垂直于底面，所以该四棱柱不是直四棱柱. 命题 \(\circled{1}\) 不成立.

对命题 \(\circled{2}\)，设侧棱方向向量为 \(\bs v\)，其在底面内的投影为 \(\bs p\). 过一对相对侧棱的截面在底面内的迹线是底面的一条对角线，设其方向向量为 \(\bs u\). 该截面由 \(\bs u\) 和 \(\bs v\) 张成. 截面垂直于底面时，截面内必有一条竖直方向，因而 \(\bs p\parallel\bs u\). 另一截面对应底面的另一条对角线，设方向向量为 \(\bs w\). 该截面垂直于底面同样说明侧棱的底面投影必须平行于其底面迹线，故 \(\bs p\parallel\bs w\). 两条对角线不平行，所以 \(\bs p=\bs0\)，即侧棱垂直于底面，命题 \(\circled{2}\) 成立.

对命题 \(\circled{3}\)，仍取底面为水平正方形，取侧棱方向向量为 \(\bs v=(1,1,1)\). 每个侧面都是由长度为 \(1\) 的底边和长度为 \(\sqrt3\) 的侧棱组成的平行四边形；侧棱与两种底边方向的内积都为 \(1\)，所以四个平行四边形的两边及夹角分别相等，四个侧面两两全等. 但 \(\bs v\) 在底面内的投影为 \((1,1,0)\ne\bs0\)，侧棱不垂直于底面，命题 \(\circled{3}\) 不成立.

对命题 \(\circled{4}\)，设底面两条对角线的方向向量为 \(\bs u\)、\(\bs w\)，侧棱方向向量为 \(\bs v\). 四条空间对角线的方向向量分别可写成 \(\bs u+\bs v\)、\(\bs u-\bs v\)、\(\bs w+\bs v\)、\(\bs w-\bs v\). 由四条对角线两两相等，特别得到
\[
\lvert\bs u+\bs v\rvert^2=\lvert\bs u-\bs v\rvert^2
\]
以及
\[
\lvert\bs w+\bs v\rvert^2=\lvert\bs w-\bs v\rvert^2
\]
展开并相减，分别得到 \(\bs u\cdot\bs v=0\) 和 \(\bs w\cdot\bs v=0\). 由于 \(\bs u\)、\(\bs w\) 不共线，它们张成底面，所以 \(\bs v\) 垂直于底面，命题 \(\circled{4}\) 成立. 因此真命题的编号为 \(2\)、\(4\).
\end{solution}

\section{解答题}
\begin{problem}
已知锐角三角形 \(ABC\) 中，\(\sin(A+B)=\frac{3}{5}\)，\(\sin(A-B)=\frac{1}{5}\). 求证：

\begin{enumerate}
\item \( \tan A = 2\tan B \)；
\item 设 \(AB = 3\)，求 \(AB\) 边上的高.
\end{enumerate}
\end{problem}

\begin{solution}
\begin{enumerate}
\item 因为 \(A\)、\(B\) 为锐角，且 \(A+B=\pi-C\in\left(\frac\pi2,\pi\right)\)，所以 \(\cos(A+B)=-\frac45\). 又 \(A-B\in\left(-\frac\pi2,\frac\pi2\right)\)，且 \(\sin(A-B)>0\)，故 \(A>B\)，从而 \(\cos(A-B)=\frac{2\sqrt6}{5}\). 由和差化积公式，
\[
2\sin A\cos B=\sin(A+B)+\sin(A-B)=\frac45
\]
\[
2\cos A\sin B=\sin(A+B)-\sin(A-B)=\frac25
\]
两式相除，并注意 \(A\)、\(B\) 为锐角，得 \(\tan A=2\tan B\).

\item 由 \(\sin(A+B)=\frac35\)、\(\cos(A+B)=-\frac45\)，得 \(\tan(A+B)=-\frac34\). 令 \(t=\tan B>0\)，由第一问 \(\tan A=2t\)，于是
\[
\frac{3t}{1-2t^2}=-\frac34
\]
即 \(2t^2-4t-1=0\). 解得 \(t=1\pm\frac{\sqrt6}{2}\)，由 \(t>0\) 得 \(t=1+\frac{\sqrt6}{2}\).

设从 \(C\) 向 \(AB\) 作垂线，垂足为 \(D\)，高为 \(CD\). 三角形为锐角三角形，故 \(D\) 在 \(AB\) 上，且
\(\tan A=\frac{CD}{AD}\)，\(\tan B=\frac{CD}{BD}\)
由 \(\tan A=2\tan B\) 得 \(BD=2AD\). 又 \(AD+BD=AB=3\)，所以 \(AD=1\)、\(BD=2\). 因此
\[
CD=BD\tan B=2+\sqrt6
\]
故 \(AB\) 边上的高为 \(2+\sqrt6\).
\end{enumerate}
\end{solution}

\begin{problem}
已知 \(8\text{ 支}\) 球队中有 \(3\text{ 支}\) 弱队，以抽签的方式将这 \(8\text{ 支}\) 球队分为 A，B 两组，每组 \(4\text{ 支}\). 求：

\begin{enumerate}
\item A，B 两组中有一组恰有 \(2\text{ 支}\) 弱队的概率；
\item A 组中至少有 \(2\text{ 支}\) 弱队的概率.
\end{enumerate}
\end{problem}

\begin{solution}
\begin{enumerate}
\item 从 \(8\) 支球队中确定 A 组的 \(4\) 支，所有分组等可能，总数为 \(\mathrm C_8^4=70\). 设 A 组有 \(k\) 支弱队，则 B 组有 \(3-k\) 支弱队. 要使两组中有一组恰有 \(2\) 支弱队，必须且只需 \(k=2\) 或 \(k=1\). 因此所求概率为
\[
\frac{\mathrm C_3^2\mathrm C_5^2+\mathrm C_3^1\mathrm C_5^3}{\mathrm C_8^4}
=\frac{30+30}{70}=\frac67
\]

\item A 组中至少有两支弱队包括恰有两支或恰有三支，所以所求概率为
\[
\frac{\mathrm C_3^2\mathrm C_5^2+\mathrm C_3^3\mathrm C_5^1}{\mathrm C_8^4}
=\frac{30+5}{70}=\frac12
\]
\end{enumerate}
\end{solution}

\begin{problem}
数列 \(\{a_{n}\}\) 的前 \(n\) 项和记为 \(S_{n}\)，已知 \(a_{1} = 1\)，\(a_{n+1} = \frac{n+2}{n} S_{n}\) （\(n = 1\)，\(2\)，\(3\)，\(\ldots\)）. 证明：

\begin{enumerate}
\item 数列 \( \left\{\frac{S_{n}}{n}\right\} \) 是等比数列；
\item \(S_{n + 1} = 4a_n\).
\end{enumerate}
\end{problem}

\begin{solution}
\begin{enumerate}
\item 由 \(a_{n+1}=S_{n+1}-S_n\) 及题设递推式，得
\[
S_{n+1}=S_n+a_{n+1}=S_n+\frac{n+2}{n}S_n=\frac{2n+2}{n}S_n
\]
两边除以 \(n+1\)，得
\[
\frac{S_{n+1}}{n+1}=2\frac{S_n}{n}
\]
又 \(S_1=a_1=1\)，所以 \(\left\{\frac{S_n}{n}\right\}\) 是首项为 \(1\)、公比为 \(2\) 的等比数列.

\item 由上题结论，\(S_n=n2^{n-1}\). 当 \(n\geqslant2\) 时，
\[
a_n=S_n-S_{n-1}=n2^{n-1}-(n-1)2^{n-2}=(n+1)2^{n-2}
\]
当 \(n=1\) 时，题设给出 \(a_1=1=(1+1)2^{-1}\)，所以该表达式对所有正整数 \(n\) 都成立. 当 \(n\geqslant2\) 时，
\[
S_{n+1}=(n+1)2^n=4(n+1)2^{n-2}=4a_n
\]
当 \(n=1\) 时，\(S_2=a_1+a_2=1+\frac{1+2}{1}S_1=4=4a_1\). 因此对任意正整数 \(n\)，都有 \(S_{n+1}=4a_n\).
\end{enumerate}
\end{solution}

\begin{problem}
如图，直三棱柱 \(ABC-A_1B_1C_1\) 中，\(\angle ACB = 90^\circ\)，\(AC = 1\)，\(CB = \sqrt{2}\)，侧棱 \(AA_1 = 1\)，侧面 \(AA_1B_1B\) 的两条对角线交点为 \(D\)，\(B_1C_1\) 的中点为 \(M\).

\begin{enumerate}
\item 证明 \(CD \perp\) 平面 \(BDM\)；
\item 求面 \(B_1BD\) 与面 \(CBD\) 所成二面角的大小.
\end{enumerate}

\bitmapfigure[max width=0.42\linewidth]{img/2004/national_paper_2_science/q20_fig1.png}
\end{problem}

\begin{solution}
\begin{enumerate}
\item 建立空间直角坐标系，取 \(C\) 为原点，令 \(CA\)、\(CB\)、\(CC_1\) 分别为三个坐标轴正向，则
\(A=(1,0,0)\)，\(B=(0,\sqrt2,0)\)，\(C_1=(0,0,1)\)
\(A_1=(1,0,1)\)，\(B_1=(0,\sqrt2,1)\).
由于 \(D\) 是矩形 \(AA_1B_1B\) 的对角线交点，且 \(M\) 是 \(B_1C_1\) 的中点，所以
\(D=\left(\frac12,\frac{\sqrt2}{2},\frac12\right)\)，\(M=\left(0,\frac{\sqrt2}{2},1\right)\).
有
\(\overrightarrow{CD}=\left(\frac12,\frac{\sqrt2}{2},\frac12\right)\)，\(\overrightarrow{BD}=\left(\frac12,-\frac{\sqrt2}{2},\frac12\right)\)，\(\overrightarrow{DM}=\left(-\frac12,0,\frac12\right)\).
直接计算得 \(\overrightarrow{CD}\cdot\overrightarrow{BD}=0\)、\(\overrightarrow{CD}\cdot\overrightarrow{DM}=0\). 由于 \(BD\)、\(DM\) 是平面 \(BDM\) 内相交的两条直线，所以 \(CD\perp\) 平面 \(BDM\).

\item 设 \(G\)、\(F\) 分别为 \(BD\)、\(BC\) 的中点. 由 \(AC\perp BC\)、\(AC=1\)、\(BC=\sqrt2\)，得
\(AB=\sqrt{AC^2+BC^2}=\sqrt3\). 由于 \(D\) 是矩形 \(AA_1B_1B\) 的中心，
\[
BD=B_1D=\frac12\sqrt{AB^2+AA_1^2}=\frac12\sqrt{3+1}=1
\]
又 \(BB_1=AA_1=1\)，所以三角形 \(BB_1D\) 为边长为 \(1\) 的正三角形. 因为
\[
\overrightarrow{GF}=-\frac12\overrightarrow{CD}
\]
所以 \(GF\perp BD\). 因而 \(B_1G\perp BD\)，且 \(\angle B_1GF\) 是所求二面角的平面角.
由坐标或勾股定理，\(B_1G=\frac{\sqrt3}{2}\)、\(GF=\frac12\)，并且 \(B_1F=\sqrt{B_1B^2+BF^2}=\sqrt{\frac32}\). 在三角形 \(B_1GF\) 中，由余弦定理，
\[
\cos\angle B_1GF=\frac{B_1G^2+GF^2-B_1F^2}{2B_1G\cdot GF}=-\frac{\sqrt3}{3}
\]
故所求二面角为 \(\arccos\left(-\frac{\sqrt3}{3}\right)=\pi-\arccos\frac{\sqrt3}{3}\).
\end{enumerate}
\end{solution}

\begin{problem}
给定抛物线 \(C: y^{2} = 4x\)，\(F\) 是 \(C\) 的焦点，过点 \(F\) 的直线 \(l\) 与 \(C\) 相交于 \(A\)，\(B\) 两点.

\begin{enumerate}
\item 设 \(l\) 的斜率为 \(1\)，求 \(\overrightarrow{OA}\) 与 \(\overrightarrow{OB}\) 夹角的大小；
\item 设 \(\overrightarrow{FB} = \lambda\overrightarrow{AF}\)，若 \(\lambda \in [4,9]\)，求 \(l\) 在 \(y\) 轴上截距的变化范围.
\end{enumerate}
\end{problem}

\begin{solution}
\begin{enumerate}
\item 抛物线 \(y^2=4x\) 的焦点为 \(F=(1,0)\). 因为直线 \(l\) 过 \(F\) 且斜率为 \(1\)，故其方程为 \(y=x-1\). 代入抛物线方程，得
\((x-1)^2=4x\)，\(x^2-6x+1=0\).
所以交点的横坐标为 \(3+2\sqrt2\)、\(3-2\sqrt2\)，相应纵坐标为 \(2+2\sqrt2\)、\(2-2\sqrt2\). 设这两个交点分别为 \(A\)、\(B\)，则
\[
\overrightarrow{OA}\cdot\overrightarrow{OB}
=(3+2\sqrt2)(3-2\sqrt2)+(2+2\sqrt2)(2-2\sqrt2)=1-4=-3
\]
并且
\(\lvert\overrightarrow{OA}\rvert^2=29+20\sqrt2\)，\(\lvert\overrightarrow{OB}\rvert^2=29-20\sqrt2\)
所以
\[
\lvert\overrightarrow{OA}\rvert^2\lvert\overrightarrow{OB}\rvert^2
=(29+20\sqrt2)(29-20\sqrt2)=29^2-800=41
\]
因此
\[
\cos\angle AOB=-\frac3{\sqrt{41}}=-\frac{3\sqrt{41}}{41}
\]
由于夹角属于 \([0,\pi]\)，故夹角为 \(\displaystyle\pi-\arccos\frac{3\sqrt{41}}{41}\).

\item 用参数表示抛物线上的点：\(P(t)=(t^2,2t)\). 设 \(A=P(t_A)\)、\(B=P(t_B)\). 两点参数为 \(u\)、\(v\) 时，弦的方程为
\[
2x-(u+v)y+2uv=0
\]
弦 \(AB\) 过焦点 \(F=(1,0)\)，代入上式得 \(2+2t_A t_B=0\)，即 \(t_A t_B=-1\). 令 \(t_B=t\)，则 \(t_A=-\frac1t\). 于是
\(\overrightarrow{FB}=(t^2-1,2t)\)，\(\overrightarrow{AF}=\left(1-\frac1{t^2},\frac2t\right)\).
由 \(\overrightarrow{FB}=\lambda\overrightarrow{AF}\) 的分量比，得 \(\lambda=t^2\). 所以 \(t=\pm\sqrt\lambda\)，其中 \(\lambda\in[4,9]\).
直线 \(l\) 过 \(F\)，其斜率为
\[
m=\frac{2t}{t^2-1}
\]
故在 \(y\) 轴上的截距为
\[
b=-m=-\frac{2t}{t^2-1}
\]
令 \(u=\lvert t\rvert=\sqrt\lambda\)，则 \(u\in[2,3]\)，且
\[
\left(\frac{2u}{u^2-1}\right)'=-\frac{2(u^2+1)}{(u^2-1)^2}<0
\]
其端点值为 \(\frac43\)、\(\frac34\). 考虑 \(t\) 的正负两种情形，截距的变化范围为
\[
\left[-\frac43,-\frac34\right]\cup\left[\frac34,\frac43\right]
\]
\end{enumerate}
\end{solution}

\begin{problem}
已知函数 \(f(x) = \ln(1 + x) - x\)，\(g(x) = x \ln x\).

\begin{enumerate}
\item 求函数 \(f(x)\) 的最大值；
\item 设 \(0 < a < b\)，证明：\(0 < g(a) + g(b) - 2g\left(\frac{a + b}{2}\right) < (b - a)\ln 2\).
\end{enumerate}
\end{problem}

\begin{solution}
\begin{enumerate}
\item 函数定义域为 \(x>-1\)，且
\[
f'(x)=\frac1{1+x}-1=-\frac{x}{1+x}
\]
当 \(-1<x<0\) 时，\(f'(x)>0\)；当 \(x>0\) 时，\(f'(x)<0\). 所以 \(f(x)\) 在 \((-1,0)\) 上递增，在 \((0,+\infty)\) 上递减，最大值在 \(x=0\) 处取得，且 \(f(0)=0\). 故最大值为 \(0\).

\item 设 \(m=\frac{a+b}{2}\)，并令 \(t=\frac{b-a}{a+b}\)，则 \(0<t<1\)，且 \(a=m(1-t)\)、\(b=m(1+t)\). 由 \(g(x)=x\ln x\)，有
\[
g(a)+g(b)-2g(m)=m\bigl((1-t)\ln(1-t)+(1+t)\ln(1+t)\bigr)
\]
令 \(H(t)=(1-t)\ln(1-t)+(1+t)\ln(1+t)\). 则
\(H'(t)=\ln\frac{1+t}{1-t}>0\)，\(H''(t)=\frac2{1-t^2}>0\).
故 \(H(t)>H(0)=0\). 又 \(H\) 在 \([0,1]\) 上严格凸，且 \(H(0)=0\)、\(H(1)=2\ln2\)，所以对 \(0<t<1\)，有 \(H(t)<2t\ln2\). 于是
\[
0<mH(t)<2mt\ln2=(b-a)\ln2
\]
所证不等式成立.
\end{enumerate}
\end{solution}
